\section{Related Work}
\label{sec:related}

\para{Memorization vs. generalization in LLMs.}
\citet{wang2025generalization} introduce \emph{distributional memorization} --- a correlation between LLM output probabilities and pretraining-corpus n-gram frequencies --- and show that for factual QA (TriviaQA), accuracy is largely explained by pretraining co-occurrences, while for math (GSM8K) and translation (WMT) it is not. Their method requires an n-gram index over the pretraining corpus, which is impractical for proprietary frontier models. \citet{lee2021deduplicating} quantify the severity of train/test overlap: $>4\%$ of standard-benchmark validation items overlap with C4, and aggressive deduplication reduces verbatim memorization by $10\times$. Unlike these approaches, we sidestep the post-hoc corpus accounting entirely by leveraging a model whose pretraining corpus is \emph{structurally} disjoint from any post-1930 concept.

\para{Contamination-controlled benchmarks.}
A complementary line of work builds benchmarks designed to be contamination-free by construction. AntiLeak-Bench~\citep{wu2024antileak} chains timestamped knowledge updates to constructively avoid overlap with a target model's training data. Daily Oracle, FreshBench, and evolveQA all exploit the train-cutoff / test-creation date gap to keep benchmark items recent. These methods must be maintained --- the cutoff moves as models do. \talkienineteen's 1930 cutoff is fixed and far enough back that \emph{any} modern benchmark is automatically contamination-free with respect to it. We argue this trades a maintenance problem for a structural assumption.

\para{Narrow-corpus generalization studies.}
\citet{li2023othello} train a GPT on Othello move sequences alone (no rules, no board state) and recover a faithful internal representation of the board state via linear probes. This is the canonical example of \emph{narrow-corpus} $\to$ \emph{controlled generalization} claim. \talkienineteen is the natural-language analogue: rather than removing the rules of a game, it removes the post-1930 world. Our work uses this structural removal to evaluate the very methodology of attribution rather than to study learned representations directly.

\para{Diagnostic probes and the limits of headline numbers.}
\citet{fu2025absencebench} show that even strong models fail at identifying \emph{omitted} content despite near-perfect needle-in-a-haystack performance. \citet{mccoy2019hans} show NLI models can rely on surface heuristics. These works share our diagnostic stance: design probes that distinguish surface success from understanding. We extend this to the attribution question --- "did the model generalize?" --- and argue that the cleanest such probe is one whose corpus disjointness is guaranteed by construction rather than verified post-hoc.

\para{The Talkie release and Holtzman's IdeaHub proposal.}
\citet{talkielm2026} release \talkienineteen alongside a matched modern twin \talkieweb and report three flagship probes: \humaneval, an \mmlu anachronism filter, and an NYT On-This-Day BpB curve. Their framing positions \talkienineteen as a tool for clean generalization claims. The IdeaHub proposal due to Holtzman (April 2026) goes further --- proposing that the sparse coverage of \talkienineteen's corpus makes generalization tests \emph{qualitatively} easier than any modern-LLM setup. To our knowledge, our paper is the first to operationalize and quantitatively test this meta-claim. The companion 4B-scale time-locked family \rankefour~\citep{ranke4b2025} provides an adjacent design point we cannot yet exploit because its weights were not publicly available at the time of writing.

\para{In-context learning of post-cutoff languages.}
The Talkie team's own \humaneval result --- a 1930-cutoff model passing Python from in-context demonstrations alone --- is the most striking generalization story currently in the literature for a cleanly time-cut model. We complement it with shorter, prompt-format-controlled ICL probes (\secref{sec:results}) that the bare 13B base model can actually answer at greedy decoding without instruction tuning, and we contrast \talkienineteen's success against its matched modern twin to isolate the cutoff effect.
